% Title page

\begin{titlepage}
\begin{flushleft}

\vspace{1cm}

\noindent\textbf{\captionfont\fontsize{28pt}{28pt}\selectfont\thesisname}

\vspace{2cm}

\noindent{\LARGE\captionfont Bogdan \textsc{Opanchuk}}

\vspace{3cm}

\noindent{\normalsize Submitted in total fulfilment of the requirements \\
of the degree of Doctor of Philosophy.}

\vfill

\includegraphics[width=4cm]{figures_precreated/swinburne_logo_and_coa.eps}

\noindent{\captionfont\scshape Centre for Atom Optics and Ultrafast Spectroscopy \\
Swinburne University of Technology \vfleuron\ Melbourne, Australia \\}

\vspace{0.5cm}

\noindent\textsc{\captionfont \today}

\end{flushleft}
\end{titlepage}


% Abstract

\cleardoublepage

\addchap*{Abstract}

% An abstract of not more than 400 words should be bound into each copy of the examinable outcome. The abstract should indicate the problem investigated, the procedures followed, the general result obtained and the major conclusions reached. It should not contain any illustrative material or tables.

A method of simulating the full quantum field dynamics of multimode multi-component Bose-Einstein condensates is developed.
The truncated Wigner representation is used to obtain a probabilistic theory.
The representation is extended to to use a functional calculus and can be applied to quantum fields in the arbitrary number of dimensions.
Detailed proofs of the corresponding theorems are given.

This method produces c-number stochastic equations which may be solved using conventional stochastic integration methods.
The technique is valid for large mode occupation numbers.
It describes spatial evolution of spinor components and properly accounts for nonlinear interactions and losses.
The method is further used to model several experiments where quantum effects are significant, and untangible for simpler approximations.
It is shown to describe leading quantum corrections correctly for effects such as phase noise, quantum squeezing, entanglement, \abbrev{epr} correlations, and interactions with engineered nonlinear reservoirs.

In conclusion, the functional Wigner representation framework developed in this thesis provides a relatively fast and accurate method of simulating dynamics of large nonlinear quantum systems, such as trapped Bose-Einstein condensates.
The methods of solving the resulting stochastic equations are easily parrallelisable and scalable, and can be implemented effectively on modern hardware.

% Acknowledgements

\cleardoublepage

\addchap*{Acknowledgements}

I would like to thank both of my supervisors, Prof.~Peter~D. Drummond and Prof.~Andrei~I. Sidorov, for many helpful discussions.

My primary supervisor, Peter Drummond, provided a fundamental outlook on the methods that I developed and used, employing both mathematical arguments and physical considerations.
He could always give a practical advice on the application of numerical methods to various problems.
And, besides that, I appreciate him encouraging me to venture into the field of parallel computing, which turned out to be perfectly suited to solve the problems we were investigating.

My secondary supervisor, Andrei Sidorov, supplied an expertise in the experimental side of my research, explaining where and how the theoretical models I used clash with reality.

I thank my fellow theorists: Rodney Polkinghorne, Dr.~Qiongyi He, Dr.~Laura Rosales-Z\'{a}rate, who gave me advice on tricky problems I was facing, and with whom I enjoyed co-authoring papers.
And in the experimental part of our research group, I thank Dr.~Michael Egorov and Valentine Ivannikov, who had a hands-on expertise on the \abbrev{bec} experiments in Swinburne, and who supplied me with data necessary to test my models. I also appreciate numerous discussions with Michael with topics ranging from quantum mechanics to programming and life in Australia.


% Declaration

\cleardoublepage

\addchap*{Declaration}

This is to certify that this thesis

\begin{itemize}
\item[\hfleuron] contains no material which has been accepted for the award to me of any other degree or diploma, except where due reference is made in the text,
\item[\hfleuron] to the best of my knowledge contains no material previously published or written by another person except where due reference is made in the text, and
\item[\hfleuron] where the work is based on joint research or publications, discloses the relative contributions of the respective workers or authors.
\end{itemize}

\vspace{2cm}

\begin{flushright}
Bogdan Opanchuk \hspace{0.5cm} \rule{4cm}{1pt} \\
\vspace{1cm}
Date \hspace{0.5cm} \rule{4cm}{1pt}
\end{flushright}
